\section{OP AMPS - performance and limitations}
\subsection{Introduction}
OP amps (Operational Amplifiers) are high-gain voltage amplifiers with differential inputs and a single-ended output. In ideal
conditions, an OP amp has infinite gain, infinite input impedance, zero output impedance and infinite bandwidth, but in real
conditions those properties are not achievable.

\subsubsection*{History}
\begin{itemize}
	\item \textbf{1953}: K2-W - first commercially available OP amp based on vacuum tubes
	\item \textbf{1947}: first transistor (germanium) at Bell Labs
	\item \textbf{1954}: first silicon transistor at Texas Instruments
	\item \textbf{1958}: first integrated circuit at Texas Instruments
	\item \textbf{1960s}: solid state op-amps
	\item \textbf{1968}: Fairchild invented the \(\mu{A741}\) (commercially available as LM741), still used today
\end{itemize}

\subsection{Ideal op amp common circuit configurations}
\subsubsection*{Circuit model of an ideal op amp}
\begin{center}
	\begin{minipage}{0.35\textwidth}
		\centering \begin{circuitikz}[american, scale=0.7, transform shape]
			
			% Disegna il triangolo esterno dell'amplificatore operazionale
			\draw[line width=1pt] (0, 2.2) -- (4, 0) -- (0, -2.2) -- cycle;
	
			% disegna i terminali di ingresso
			\draw (-0.5, 1) node[left] (vplus) {$V^+$} to[short, o-] ++(0.85, 0)
				to[R, l=$R_{\mathrm{in}}$] ++(0, -2)
				to[short, -o] ++(-0.85, 0) node[left] (vminus) {$V^-$};
	
			% disegna la tensione di ingresso
			\draw[-{Latex}]
				(vminus) to[out=-250, in=250] 
				node[midway, fill=white, inner sep=2pt] {$V_{in}$} 
				(vplus);
	
			% generatore di tensione e resistenza
			\draw (4.5, 0) node[right] (vout) {$V_{out}$}
				to[short, o-] ++(-1, 0)
				to[R, l_=$R_{out}$] ++(-1.8, 0)
				to[cV, name=cgen] ++(0, -1.2)
				to[short] ++(0, -0.2) node[ground] {};
			\node at (cgen) [xshift=0.9cm, yshift=-0.4cm, fill=white, inner sep=1pt] {\(A \! \cdot \! V_{in}\)}; 
			
			% pin di alimentazione
			\draw (1.2, 1.55) to[short, -o] ++(0, 0.7) node[above] {$V_{DD+}$};
			\draw (1.2, -1.55) to[short, -o] ++(0, -0.7) node[below] {$V_{SS-}$};
	
		\end{circuitikz}
	\end{minipage}
	\begin{minipage}{0.6\textwidth}
		\begin{itemize}
			\item \(A\): open-loop gain of the op amp
			\item \(R_{in}\): input resistance of the op amp
			\item \(R_{out}\): output resistance of the op amp
			\item \(V_{in}\): differential input voltage (\(V_{in} = V^+ - V^-\))
			\item \(V_{out}\): output voltage
			\item \(V_{DD+}\) and \(V_{SS-}\): power supply voltages for the op amp
		\end{itemize}
	\end{minipage}
\end{center}

\subsubsection*{Inverting amplifier}
\begin{center}
	\begin{minipage}{0.35\textwidth}
		\centering \begin{circuitikz}[american, scale=0.8, transform shape]
			\node[op amp] (opamp) at (0,0) {};
			\draw (opamp.+) to[short] ++(0,-0.5) node[ground] {};
			\draw (opamp.-) to[R, l_=$R_{in}$] ++(-2,0) 
				node[left] {$V_{in}$}
				to[short, o-] ++(0,0);
			
			\draw (opamp.-) to[short, *-] ++(0, 1) coordinate (fb_up)
				to[R, l=$R_{f}$] (fb_up -| opamp.out)
				to[short, -*] (opamp.out);

			\draw (opamp.out) to[short, -o] ++(0.5,0) node[right] {$V_{out}$};
		\end{circuitikz}
	\end{minipage}
	\begin{minipage}{0.63\textwidth}
		Differential input voltage \(V_{in} = V^- - V^+\) must be zero to have a finite output voltage. For this reason,
		\(V^- = V^+ = 0V\) and the inverting input is called a virtual ground potential node because it is at ground
		potential but not directly connected to ground.
		
		The resistor series \(R_{in}\) and \(R_f\) acts as a voltage divider and the output voltage \(V_{out}\) is proportional
		to the input voltage \(V_{in}\) by a factor \(-R_f / R_{in}\):
		\[V_{out} = -\frac{R_f}{R_{in}} V_{in} \qquad H(s) = \frac{V_{out}}{V_{in}} = -\frac{R_f}{R_{in}}\]

		The output signal is out of phase (inverted) with the input signal.
	\end{minipage}
\end{center}

\subsubsection*{Non-inverting amplifier}
\begin{center}
	\begin{minipage}{0.35\textwidth}
		\centering \begin{circuitikz}[american, scale=0.8, transform shape]
			\node[op amp] (opamp) at (0,0) {};
			\draw (opamp.+) to[short, -o] ++(0, 0) node[left] {\(V_{in}\)};
			\draw (opamp.-) to[R, l_=$R_{in}$] ++(-2,0) node[ground] {};
			\draw (opamp.out) to[short, -o] ++(0.5,0) node[right] {$V_{out}$};

			\draw (opamp.-) to[short, *-] ++(0, 1) coordinate (fb_up)
				to[R, l=$R_{f}$] (fb_up -| opamp.out)
				to[short, -*] (opamp.out);
		\end{circuitikz}
	\end{minipage}
	\begin{minipage}{0.63\textwidth}
		To have a finite output voltage and a zero input differential voltage, \(V^+ = V^-\).
		
		The resistor series \(R_{in}\) and \(R_f\) acts as a voltage divider and the output voltage \(V_{out}\)
		is proportional to the input voltage \(V_{in}\) by a factor \(1 + R_f / R_{in}\):
		\[V_{out} = \left(1 + \frac{R_f}{R_{in}}\right) V_{in} \qquad H(s) = \frac{V_{out}}{V_{in}} = 1 + \frac{R_f}{R_{in}}\]

		The output signal is in phase (non-inverted) with the input signal.
	\end{minipage}
\end{center}

\subsubsection*{Voltage follower - buffer}
\begin{center}
	\begin{minipage}{0.35\textwidth}
		\centering \begin{circuitikz}[american, scale=0.9, transform shape]
			\node[op amp, xscale=0.6, yscale=0.6] (opamp) at (0,0) {};
			\draw (opamp.+) to[short, -o] ++(0, 0) node[left] {\(V_{in}\)};
			\draw (opamp.out) to[short, -o] ++(0.5,0) node[right] {$V_{out}$};

			\draw (opamp.-) to[short] ++(0, 0.6) coordinate (fb_up)
				to[short] (fb_up -| opamp.out)
				to[short, -*] (opamp.out);
		\end{circuitikz}
	\end{minipage}
	\begin{minipage}{0.63\textwidth}
		The voltage follower is a specific configuration of the non-inverting amplifier with no input and feedback resistors.
		The output signal is ideally an exact copy of the input signal from the non inverting input.
	\end{minipage}
\end{center}

\subsection{Ideal assumptions}
\begin{itemize}
	\item \textbf{infinite open-loop gain} \(A \to \infty\): the output voltage \(V_{out}\) can be arbitrarily large for any
	non-zero input voltage \(V_{in}\)
	\item \textbf{infinite input resistance} \(R_{in} \to \infty\): no current flows into the input terminals (zero input-bias
	current \(I_B = 0\)) and consequently the difference between the two input-bias currents is zero (zero input-offset current
	\(I_{OS} = 0\))
	\item \textbf{zero input offset voltage}: the output voltage is zero when the input voltage is zero
	\item \textbf{infinite common-mode rejection ratio (CMRR)}: the output voltage is only affected by the difference between
	the input voltages \(V^+ - V^-\) and not by the common-mode voltage \((V^+ + V^-)/2\)
	\item \textbf{zero output resistance} \(R_{out} \to 0\): the output voltage \(V_{out}\) is not affected by the load connected
	\item \textbf{infinite output voltage range}: the output voltage \(V_{out}\) is not limited by the power supply voltages
	and the op amp does not saturate
	\item \textbf{power supply rejection}: the output voltage is not affected by power supply voltage variations
	\item \textbf{infinite output current}: the op amp can provide any amount of current to the load
	\item \textbf{infinite open-loop bandwidth}: any frequency can be amplified without attenuation
	\item \textbf{infinite slew rate}: the output voltage can change instantaneously in response to changes in the input differential voltage
\end{itemize}

\subsection{Non idealities of op amps}
\subsubsection*{Input offset voltage}
When both inputs are connected to ground, the output has a small offset voltage \(V_0\) that keeps the output at some DC voltage
level instead of 0 V. This output voltage can be measured by connecting the op amp as a voltage follower with the non inverting
input connected to ground.

It is possible to create a model that include this effect by adding an ideal voltage generator in series to one of the two inputs
to generates a fixed input-voltage offset \(V_{OS}\). The \(V_{OS}\) is the equivalent differential voltage that must be provided
to an ideal op amp to obtain the output voltage offset \(V_0\).

Some ICs have specific pins to add a potentiometer that can be used to compensate the input offset voltage.

\begin{center}
	\begin{minipage}{0.45\textwidth}
		\centering \begin{circuitikz}[american, scale=0.8, transform shape]
			\node[op amp] (opamp) at (0,0) {};
			\draw (opamp.+) to[V, l=$V_{\mathrm{OS}}$] ++(-1,0) to[short, -o] ++(-0.5, 0) node[left] {\(V_{in}\)};
			\draw (opamp.-) to[short] ++(-1.5,0) to[R, l_=$R_{in}$] ++(-2,0) node[ground] {};
			\draw (opamp.out) to[short, -o] ++(0.5,0) node[right] {$V_{out}$};

			\draw (opamp.-) to[short] ++(-1.5,0) to[short, *-] ++(0, 1) coordinate (fb_up)
				to[R, l=$R_{f}$] (fb_up -| opamp.out)
				to[short, -*] (opamp.out);
			
			\draw[color=red] (-2.4,1.1) -- ++(3.4,0) -- ++(0,-2.7) -- ++(-3.4,0) -- cycle;
		\end{circuitikz} \\[5pt]
		\footnotesize Ideal op amp with zero output offset voltage \(V_0\) with non-ideal input offset voltage \(V_{OS}\)
		correction highlighted in the red rectangle.
	\end{minipage}
	\hfill
	\begin{minipage}{0.27\textwidth}
		\vspace{18pt}
		\centering \begin{circuitikz}[american, scale=0.8, transform shape]
			\node[op amp] (opamp) at (0,0) {};
			\draw (opamp.+) to[short] ++(0, 0) node[ground] {};
			\draw (opamp.out) to[short, -o] ++(0.5,0) node[right] {$V_{out}$};

			\draw (opamp.-) to[short] ++(0, 0.8) coordinate (fb_up)
				to[short] (fb_up -| opamp.out)
				to[short, -*] (opamp.out);
		\end{circuitikz} \\[10pt]
		\footnotesize Voltage follower configuration in which the output voltage is equal to the input offset voltage \(V_{OS}\).
	\end{minipage}
	\hfill
	\begin{minipage}{0.25\textwidth}
		\centering \begin{circuitikz}[american, scale=0.8, transform shape]
			\node[op amp] (opamp) at (0,0) {};
			\draw (opamp.+) to[short, -o] ++(0,0) node[left] {$V^+$};
			\draw (opamp.-) to[short, -o] ++(0,0) node[left] {$V^-$};
			\draw (opamp.out) to[short, -o] ++(0,0) node[right] {$V_{out}$};
			
			\draw[scale=0.5, transform shape] (-0.9, 1.5) -- ++(0, 0.7) to[R] ++(2, 0) -- ++(0, -1.9);

			\draw[<-, >=stealth] (0, 1.25) -- (0, 1.8) node[above] {\(V_{CC}\) or \(V_{EE}\)};
		\end{circuitikz} \\[5pt]
		\footnotesize Potentiometer for input offset voltage compensation
	\end{minipage}
\end{center}

Numerical example: if a non-inverting amplifier with \(R_f = 99 \; k\Omega\) and \(R_{in} = 1.2 \; k\Omega\) has an input offset
voltage of \(\left|V_{OS}\right| \leq 3 \; mV\), the output offset voltage \(V_0\) is:
\[\left|V_0\right| = \left|V_{OS}\right| \cdot \left(1 + \frac{R_f}{R_{in}}\right) = 3 \; mV \cdot \left(1 + \frac{99 \; k\Omega}{1.2 \; k\Omega}\right) = 0.25 \; V \quad \longrightarrow \quad -0.25 \; V < V_0 < 0.25 \; V\]

\subsubsection*{Input bias and offset currents}
Due to finite input resistance, there are small input bias currents \(I_{B_+}\) and \(I_{B_-}\) that flow into the input terminals
of the op amp. This happens because the input terminals are connected to the base or gate of internal transistors (BJT or MOSFET)
that have a non-infinite resistance. The difference between the two input bias currents is called input offset current \(I_{OS}\).

It is possible to create a model that include this effect by adding two ideal current generators in parallel to the input terminals
to simulate the input bias currents \(I_{B_+}\) and \(I_{B_-}\).

These currents cause voltage drops across the resistors \(R_{in}\) and \(R_f\) and consequently generate a small input
voltage offsets that will be amplified to a considerable output voltage offset. To mitigate this effect, it is recommended to
add a compensation resistor \(R_{B}\) between the non-inverting input and ground with a value equal to the parallel of \(R_{in}\)
and \(R_f\). Using this value, the output voltage offset depends only by the input offset current \(I_{OS}\) that is usually
5-10 times smaller than the individual input bias currents \(I_{B_+}\) and \(I_{B_-}\).
\[V_{out} = V_{out,1} + V_{out,2} = I_{B_-} \cdot R_f - I_{B_+} \cdot R_{B} \cdot \left(1 + \frac{R_f}{R_{in}} \right) \qquad \text{using superposition}\]
\[\text{using} \;\; R_B = R_f \parallel R_{in} = \frac{R_{in} \cdot R_f}{R_{in} + R_f} \quad \rightarrow \quad V_{out} = - (I_{B_+} - I_{B_-}) \cdot R_f = -I_{OS} \cdot R_f\]

\begin{center}
	\begin{minipage}{0.55\textwidth}
		\centering \begin{circuitikz}[american, scale=0.8, transform shape]
			\node[op amp] (opamp) at (0,0) {};
			\draw (opamp.+) to[short, -o] ++(-1.3, 0) node[above] {\(V_{in}\)};
			\draw (opamp.-) to[short] ++(-3.5,0) to[R, l_=$R_{in}$] ++(-2,0) node[ground] {};
			\draw (opamp.out) to[short, -o] ++(0.5,0) node[right] {$V_{out}$};

			\draw (opamp.-) to[short] ++(-3.5,0) to[short, *-] ++(0, 1) coordinate (fb_up)
				to[R, l=$R_{f}$] (fb_up -| opamp.out)
				to[short, -*] (opamp.out);
			
			\draw ($(opamp.+)+(-0.1,0)$) to[isource, name=ibplus] ++(0,-1.5) ++(0,0.18) node[ground] {};
			\draw ($(opamp.+)+(-0.9,0)$) to[R, l_=$R_B$] ++(0,-1.5) ++(0,0.18) node[ground] {};

			\draw ($(opamp.-)+(-2.5,0)$) to[isource, name=ibminus] ++(0,-1.5) ++(0,0.35) node[ground] {};
			
			\node at (ibplus) [xshift=0.6cm, yshift=-0.5cm, inner sep=1pt] {\(I_{B_+}\)}; 
			\node at (ibminus) [xshift=-0.4cm, yshift=-0.5cm, inner sep=1pt] {\(I_{B_-}\)}; 

			\draw[color=red] (-4.5,1.1) -- ++(5.5,0) -- ++(0,-3.6) -- ++(-5.5,0) -- cycle;
		\end{circuitikz}
	\end{minipage}
	\begin{minipage}{0.4\textwidth}
		Op amp model with current generators representing input bias currents \(I_{B_+}\) and \(I_{B_-}\) and compensation resistor
		\(R_B\) highlighted in the red rectangle.
	\end{minipage}
\end{center}

\subsubsection*{Input bias currents in integrators}
An integrator is a circuit configuration that integrates an input signal by storing charge on a capacitor \(C\) in the
feedback loop. The output voltage is proportional to the integral of the input signal. The capacitor has a normally open
switch in parallel to discharge the capacitor and reset the output voltage to the initial value.

\begin{center}
	\begin{minipage}{0.4\textwidth}
		\centering \begin{circuitikz}[american, scale=0.8, transform shape]
			\node[op amp] (opamp) at (0,0) {};
	
			\draw (opamp.out) to[short, -o] ++(0.5,0) node[right] {$V_{out}$};

			\draw (opamp.-) to[short] ++(-1,0) to[short, *-] ++(0, 1.3) coordinate (fb_up)
				to[C, name=c] (fb_up -| opamp.out)
				to[short, -*] (opamp.out);
			\node at (c) [xshift=0.4cm, yshift=-0.3cm, inner sep=1pt] {\(C\)}; 
			
			\draw (opamp.-) to[short] ++(-1,0) to[short, *-] ++(0, 2) coordinate (fb_up2)
				to[switch] (fb_up2 -| opamp.out)
				to[short] (opamp.out);
			
			\draw (opamp.+) -- ++(-0.2,0) to[vsource, name=ibplus] ++(0,-1.5) ++(0,0.35) node[ground] {};
			\draw (opamp.-) -- ++(-1.7,0) to[isource, l=\(I_{B_-}\)] ++(0,-1.5) ++(0,0.18) node[ground] {};
			\draw (opamp.-) -- ++(-2.5,0) to[R, l_=$R$] ++(0,-1.5) ++(0,0.18) node[ground] {}; 
			\node at (ibplus) [xshift=0.75cm, yshift=-0.4cm, inner sep=1pt] {\(V_{OS}\)};
			%\node at (ibminus) [xshift=-0.4cm, yshift=-0.5cm, inner sep=1pt] {\(I_{B_-}\)};
		\end{circuitikz}
	\end{minipage}
	\begin{minipage}{0.55\textwidth}
		NOTES:
		\begin{itemize}
			\item \(V_{OS}\) models the input offset voltage of real op amps
			\item \(I_{B_-}\) models the input bias current of real op amps
			\item useful formulas: \(V_C = Q_C / C\), \(\;\; Q_C = I_C \cdot t\)
		\end{itemize}
	\end{minipage}
\end{center}

\noindent
Analyzing the circuit behavior in the two possible states of the switch:
\begin{itemize}
	\item[1.] when the switch is closed, op amp acts as a voltage follower and \(V_{out} = V_{OS}\)
	\item[2.] when the switch is open, the output voltage \(V_{out}\) is:
	\[V_{out}(t) = V_{start} + V_{out,1} + V_{out,2} = V_{OS} + \frac{V_{OS}}{RC} t + \frac{I_{B}}{C} t\]
	\begin{itemize}[topsep=0pt]
		\item the input offset voltage \(V_{OS}\) that generates a current \(I_{OS} = V_{OS} / R\) through the resistor \(R\)
		that charge the capacitor \(C\) with a voltage contribution \(V_{out,1} = t \cdot V_{OS} / RC \)
		\item the input bias current \(I_{B}\) charge the capacitor \(C\) with a voltage contribution \(V_{out,2} = t \cdot I_{B} / C\)
		\item output voltage increases linearly with time and the slope is \(V_{OS} / RC + I_B / C\)
	\end{itemize}
\end{itemize}

\subsubsection*{Output voltage swing}
The output voltage swing is the range of output voltage that an op amp can provide without saturating. It is limited by 1-2 volts
below the power supply span. The saturation voltages \(L_+\) and \(L_-\) are asymmetrical and depend on the load connected to the
output and on the output current.

\subsubsection*{Output current limits}
The output current of an op amp is limited by short-circuit and power dissipation protection inside the op amp.

Sometimes datasheets specify only the output swing based on the load resistance \(R_L\), but not the maximum output current
\(I_{out, max}\). That value can be estimated by the following formula:
\[I_{out, max} = \frac{V_{out, max}}{R_L}\]

The output current is given by the sum of the load current \(I_L\) and the feedback current \(I_f\). Usually \(R_f\) is much larger
than \(R_L\) and the feedback current \(I_f\) is negligible compared to the load current \(I_L\). In general, for an inverting
amplifier the total output current \(I_{out}\) is:
\[I_{out} = I_L + I_f = \frac{V_{out}}{R_L} + \frac{V_{out}}{R_f + R_{in}} = V_{out} / R_{eq} \qquad \text{with} \quad R_{eq} = R_L \parallel (R_f + R_{in})\]

\begin{center}
	\begin{circuitikz}[american, scale=0.8, transform shape]
		\node[op amp] (opamp) at (0,0) {};
		\draw (opamp.+) to[short, -o] ++(0, 0) node[left] {\(V_{in}\)};
		\draw (opamp.-) to[R, l_=$R_{in}$] ++(-2,0) node[ground] {};
		\draw (opamp.out) -- ++(1, 0) to[R, l=$R_L$, i>^=$I_L$] ++(3,0) node[ground] {};
		\draw (opamp.out) ++(-0.3, 0) to[short, i=$I_{out}$] ++(0.7, 0);

		\draw (opamp.-) to[short, *-] ++(0, 1) coordinate (fb_up)
			to[R, l=$R_{f}$, ] (fb_up -| opamp.out)
			to[short] ++(0.7, 0) coordinate (fb_out)
			to[short, -*, i<=$I_f$] (fb_out |- opamp.out)
			to[short, -o] ++(0,-0.5) node[right] {\(V_{out}\)};
	\end{circuitikz}
\end{center}

\noindent
To correctly design the circuit and choose the correct feedback resistor \(R_f\) and input resistor \(R_{in}\) of an inverting
amplifier based on the load resistance \(R_L\), the following steps can be followed:
\begin{itemize}
	\item[1.] calculate the minimum equivalent resistance \(R_{eq} = R_L \parallel R_f\) based on \(V_{out, max}\) and \(I_{out, max}\)
	\item[2.] fix the load resistance based on the application requirements and calculate the minimum feedback resistor \(R_f\)
	(it is recommended to choose \(R_f\) slightly bigger than the minimum value)
	\item[3.] choose the input resistor \(R_{in}\) based on the desired gain of the amplifier and the value of \(R_f\) chosen
	in the previous step
	\item[4.] calculate the maximum output current \(I_{out}\) for the chosen resistor values and verify the requirements
\end{itemize}
It is recommended to choose resistor values between input impedance (\(1 \; M\Omega \; - \; 10 \; M\Omega\)) and output impedance
(\(10 \; \Omega \; - \; 100 \; \Omega\)) of the op amp. Usually, \(1 \; k\Omega \; - \; 100 \; k\Omega\) should be a good choice.

\subsubsection*{Slew rate and full-power bandwidth}
Slew rate is the maximum velocity with which the output voltage can change in response to a step input voltage. It is measured
in volts per microsecond (V/µs), usually in range \(0.1 \; V/\mu s \; - \; 10 \; V/\mu s\). If the input slew rate is higher than the op amp slew rate, the output will be clipped, filtered
or distorted.

The full-power bandwidth is the highest frequency at which a full-scale signal can be developed at the output without distortion.
It is related to the slew rate \(SR\) and the maximum output voltage \(V_{out, max}\) by the following formula:
\[f_{FPBW} = \frac{SR}{2 \pi \, V_{out, max}}\]
\[\text{\footnotesize sinusoidal output signal: } V_{out} = V_M \sin (\omega t) \quad \rightarrow \quad \text{\footnotesize signal slew rate: } \left. \frac{d V_{out}}{d t} \right|_{max} = \left. V_M \omega \cos (\omega t) \right|_{max} = V_M \omega\]
\[\text{\footnotesize signal SR should be less than the op amp SR: } \quad V_M \omega < SR \quad \rightarrow \quad V_M < \frac{SR}{\omega}\]
\[\text{\footnotesize obtaining the full-power bandwidth: } \quad \omega = 2 \pi f_M < \frac{SR}{V_M} \quad \rightarrow \quad f_M < \frac{SR}{2 \pi V_M}\]
